% ==========================================================================
% Project Design Justifications
%
% Captures the reasoning behind key architectural and design decisions.
% Use this as a lookup when writing the dissertation methodology chapter.
%
% Compile with: pdflatex justifications.tex
% ==========================================================================

\documentclass[11pt,a4paper]{article}
\usepackage[utf8]{inputenc}
\usepackage[margin=2.5cm]{geometry}
\usepackage{natbib}
\usepackage{hyperref}
\usepackage{enumitem}
\usepackage{booktabs}
\usepackage{longtable}

\title{Design Justifications\\
\large Uncertainty-Conditioned RL with Evidential Policy Networks}
\author{Antonio Galdes}
\date{Last updated: 8 April 2026}

\begin{document}
\maketitle

% ------------------------------------------------------------------
\section{Localisation Strategy: RTK-GNSS + IMU}
% ------------------------------------------------------------------

\subsection{Why not Cartographer / scan-matching}

The initial project iteration used Google Cartographer for 2D LiDAR
scan-matching against perimeter cones and parked car bodies. This was
abandoned for three reasons:

\begin{enumerate}
  \item \textbf{Cone occlusion in full lots}: When bay occupancy is high,
        parked cars occlude the perimeter cones that Cartographer relied on
        for scan-matching features. Localisation degraded catastrophically --
        the exact opposite of the desired ``graceful degradation'' behaviour.
  \item \textbf{Pipeline complexity}: The Cartographer stack added massive
        infrastructure: pbstream mapping workflow, 4 Lua configuration files,
        a custom \texttt{TfToOdomNode} for dynamic covariance extraction, an
        \texttt{ImuFrameRelay} node for frame convention translation, and
        cross-distro DDS workarounds (Humble/Jazzy). None of this contributed
        to the thesis.
  \item \textbf{Thesis contribution clarity}: The dissertation contributes
        uncertainty-conditioned RL, not a localisation method. RTK-GNSS is a
        well-understood, commercially deployed technology that produces the
        covariance variation the policy needs without requiring novel
        localisation research.
\end{enumerate}

\subsection{Why RTK-GNSS}

RTK-GNSS provides cm-level accuracy under nominal conditions (RTK fixed
integer solution), with well-characterised degradation modes:

\begin{center}
\begin{tabular}{llcl}
\toprule
Fix State & Horizontal Error (95\%) & Fits 2.5\,m bay? & Policy should... \\
\midrule
RTK fixed     & 1--3\,cm   & Yes, centred & Park normally \\
RTK float     & 20--50\,cm & Yes, tight   & Slow down, careful approach \\
DGNSS         & 0.5--1.5\,m & No          & Abort or creep \\
Standalone    & 2--5\,m    & No, wrong bay & Safety handoff \\
\bottomrule
\end{tabular}
\end{center}

A base station + rover costs approximately 600\,GBP total and is trivially
deployable at an outdoor parking facility. This is realistic: commercial
autonomous parking infrastructure (Bosch, Valeo) assumes instrumented lots.

\subsection{Outdoor-only limitation}

RTK-GNSS requires sky visibility. This system therefore targets \textbf{outdoor
parking lots only}. This is stated upfront as a scope limitation in the
dissertation. Indoor/covered parking would require a different localisation
approach (e.g.\ UWB beacons, visual SLAM) but the uncertainty-conditioned
policy architecture would transfer unchanged -- only the covariance source
changes.

\subsection{RTK fix states as natural uncertainty source}

The RTK fix-state tiers create natural discrete uncertainty regimes that
the policy must learn to handle. During training, a tier is sampled
per-episode from \texttt{configs/gnss\_noise\_profiles.yaml} with weights
favouring conditions where parking IS possible (70\% rtk\_fixed + float)
so the agent learns to park, but including enough degraded episodes (30\%)
to learn when NOT to try.

% ------------------------------------------------------------------
\section{Why Uncertainty Conditioning Matters}
% ------------------------------------------------------------------

\subsection{The seatbelt argument}

Under nominal conditions (RTK fixed, familiar layout), all baselines perform
identically. The uncertainty-conditioned system adds no visible overhead. It
earns its value in the 20--30\% of episodes where conditions degrade -- like a
seatbelt that is invisible until the crash. The ablation quantifies exactly
how much each uncertainty layer contributes to safety under degradation, while
confirming zero performance cost under nominal operation.

\subsection{Safety-critical systems defined by failure modes}

ISO 26262 requires automotive systems to detect and respond to faults, not
just operate correctly under nominal conditions. A parking system that works
perfectly 80\% of the time but crashes into other vehicles the remaining 20\%
is not deployable. The uncertainty-conditioned policy detects degradation and
triggers safety handoffs, converting potential collisions into safe stops.

\subsection{Quantified via ablation}

The key metric is \textbf{collision rate under degraded GNSS}, not success
rate under nominal conditions. The 2$\times$2 ablation study (Section 4)
isolates the contribution of each uncertainty layer:
\begin{itemize}
  \item \texttt{vanilla\_ppo}: no uncertainty awareness -- crashes under
        degraded GNSS.
  \item \texttt{input\_uncertainty}: sees EKF covariance but cannot express
        ``I don't know what to do'' -- either acts or doesn't.
  \item \texttt{output\_uncertainty}: can express uncertainty but doesn't
        know WHY (localisation vs.\ novelty).
  \item \texttt{full\_method}: both layers -- detects high covariance AND
        high epistemic uncertainty, triggers graduated response.
\end{itemize}

\subsection{Cost vs.\ risk}

The evidential actor adds negligible inference overhead (single forward pass,
4 extra NIG parameters per action dimension). The cost of NOT having
uncertainty awareness is vehicle collision in a shared parking facility.

% ------------------------------------------------------------------
\section{Two-Layer Uncertainty Architecture}
% ------------------------------------------------------------------

\subsection{Layer 1: Localisation uncertainty (EKF covariance)}

The EKF posterior covariance -- 9 elements including cross-covariance terms
(\texttt{std\_x, std\_y, std\_yaw, cov\_xx, cov\_yy, cov\_yawyaw, cov\_xy,
cov\_xyaw, cov\_yyaw}) -- appears as explicit dimensions of the RL
observation vector (indices 3--11). This captures ``how sure am I about
where I am?'' directly from sensor degradation.

Rather than compressing localisation uncertainty into scalar reward signals
(as in Active SLAM formulations) or maintaining internal belief states
(POMDP approaches), we expose the full EKF posterior. The policy learns
to interpret all 9 dimensions, including cross-covariance terms that
indicate correlated position/heading uncertainty.

\subsection{Layer 2: Action uncertainty (evidential NIG on the actor)}

The evidential actor outputs Normal-Inverse-Gamma (NIG) parameters
$(\gamma, \nu, \alpha, \beta)$ per action dimension, yielding epistemic
and aleatoric uncertainty in a single forward pass. This captures
``how sure am I about what to do?''

Unlike EPPO (Akgul et al., 2025), which applies evidential deep learning
to the critic's value function, we place the NIG evidential head on the
actor network itself. This yields per-action epistemic and aleatoric
uncertainty at inference time, enabling direct uncertainty-aware action
modulation without requiring critic evaluation.

\subsection{Why both layers are needed}

Each layer responds to a distinct failure mode:

\begin{center}
\begin{tabular}{lcccp{4.5cm}}
\toprule
Scenario & Loc.\ uncert. & Epistemic & Aleatoric & Correct behaviour \\
\midrule
RTK fixed, familiar bay & Low & Low & Low & Park normally \\
RTK fixed, novel bay angle & Low & \textbf{High} & Low &
  Park successfully; epistemic flags novelty (ODD boundary detector) \\
RTK float, familiar bay & Medium & Low & \textbf{High} &
  Drive slowly (noisy outcomes) \\
RTK lost, any bay & \textbf{High} & \textbf{High} & \textbf{High} &
  Safety handoff \\
\bottomrule
\end{tabular}
\end{center}

\subsection{OOD calibration argument}

Elevated epistemic uncertainty on unseen layouts is proof of correct
calibration, not proof of behaviour change. The evidential actor correctly
detects it is outside its training distribution even when the task happens
to be manageable. The value materialises when OOD + degraded localisation
stack: the system that correctly flags ``I haven't seen this before'' when
things are easy is the same system that saves the car when things get hard.

% ------------------------------------------------------------------
\section{Ablation Study Design}
% ------------------------------------------------------------------

\subsection{2$\times$2 factorial}

\begin{center}
\begin{tabular}{l|cc}
 & EKF covariance in obs? No & Yes \\
\hline
Evidential actor? No  & \texttt{vanilla\_ppo} (11d) & \texttt{input\_uncertainty} (20d) \\
Evidential actor? Yes & \texttt{output\_uncert} (11d) & \texttt{full\_method} (20d) \\
\end{tabular}
\end{center}

\subsection{Why these 4 baselines are the minimum sufficient set}

Each cell isolates one factor while controlling the other. The 2$\times$2
design answers two questions independently: (1) does seeing localisation
uncertainty help? (2) does expressing action uncertainty help? And crucially:
(3) does combining both produce a synergy beyond either alone?

\subsection{Key metrics}

\begin{itemize}
  \item \textbf{Success rate (nominal)}: can it park when RTK is good?
        All 4 should be similar.
  \item \textbf{Success rate (degraded)}: can it park when RTK drops?
        \texttt{full\_method} should degrade gracefully.
  \item \textbf{Collision rate (degraded)}: does it crash when uncertain?
        \texttt{full\_method} should approach 0; \texttt{vanilla\_ppo} crashes.
  \item \textbf{False attempt rate}: how often does it try to park when
        it shouldn't? \texttt{full\_method} should be lowest.
  \item \textbf{Graceful degradation curve}: success vs.\ GNSS noise level.
        \texttt{full\_method} drops last.
  \item \textbf{Epistemic uncertainty (OOD layout)}: does it know when it
        is in an unseen lot? Evidential baselines should show higher values.
  \item \textbf{Aleatoric uncertainty (RTK float)}: does it detect noisy
        outcomes? Evidential baselines should show higher values.
\end{itemize}

% ------------------------------------------------------------------
\section{Sensor Architecture}
% ------------------------------------------------------------------

\begin{center}
\begin{tabular}{llll}
\toprule
Sensor & Role & Feeds into & Real-world cost \\
\midrule
RTK-GNSS (base + rover) & Position fixes & EKF (localisation) & $\sim$600\,GBP \\
IMU & Orientation, yaw rate & EKF (prediction step) & $\sim$50\,GBP \\
2D LiDAR & Obstacle clearance & Obs indices 15--19 directly & $\sim$500--1500\,GBP \\
\bottomrule
\end{tabular}
\end{center}

Sensor roles are cleanly separated: RTK-GNSS + IMU = localisation
(``where am I?''). LiDAR = obstacle detection (``what's near me?'').
Each sensor has exactly one job. LiDAR does not participate in localisation.

Wheel odometry excluded: CARLA bridge publishes ground-truth physics state
as odometry (zero noise), which suppresses covariance variation.

% ------------------------------------------------------------------
\section{Novelty and Positioning in Literature}
% ------------------------------------------------------------------

The system's novelty rests on three independently novel pillars whose
combination is unprecedented:

\subsection{Pillar 1: Evidential actor}

NIG on the actor of a continuous-action PPO policy. No prior work places
EDL on the actor; all existing EDL-in-RL is on critics/Q-functions.
Closest: EPPO (Akgul et al., 2025) applies EDL to critic V(s) only;
actor remains standard Gaussian.

\subsection{Pillar 2: Localisation covariance as state}

Raw EKF/GNSS covariance (9 elements including cross-covariance) as RL
observation dimensions for vehicle control. No prior work for ego-vehicle
localisation. Closest: Cai et al.\ (2024) feeds pedestrian prediction
covariance into state (different uncertainty source and purpose).

\subsection{Pillar 3: Dual-layer uncertainty for parking}

The ``where am I uncertain?'' / ``what should I do?'' decomposition
applied to autonomous parking. No prior work on uncertainty-conditioned
parking RL of any kind. A 2025 MDPI Sensors survey identifies uncertainty
estimation as ``critical'' for parking safety but notes it exists only at
perception level, not control.

\subsection{Key competitors}

\begin{center}
\begin{longtable}{lclp{5.5cm}}
\toprule
Paper & Year & Threat & Why different \\
\midrule
EPPO (Akgul et al.) & 2025 & Medium & EDL on critic V(s), not actor. MuJoCo, not driving. \\
CEQR-DQN (Stutts et al.) & 2024 & Low & EDL on Q-network. Discrete actions (MinAtar). \\
DRO-EDL-MPC (Ham \& Ahn) & 2025 & Low & EDL on perception, not policy. MPC, not RL. \\
UNICON (IEEE TNNLS) & 2022 & Low-Med & Ensemble disagreement conditions policy. Not evidential. MuJoCo. \\
Cai et al. & 2024 & Low-Med & Pedestrian prediction covariance in state. Crowd nav, not parking. \\
Active SLAM DRL & 2026 & Low & Scalar SLAM covariance in reward. Exploration, not parking. \\
\bottomrule
\end{longtable}
\end{center}

\subsection{RTK-GNSS covariance as RL input}

Completely novel. No paper feeds GNSS-derived covariance or quality metrics
into an RL policy for vehicle control. The RTK fix-state tiers (fixed
$\sim$2\,cm, float $\sim$30\,cm, standalone $\sim$2--5\,m) create natural
discrete uncertainty regimes with no precedent in the RL literature.

\end{document}
